% =========================================================
\chapter{Lecture 3 : Access control basics and DAC}
% =========================================================

\begin{rembox}[title=In a nutshell]
\textbf{What is this lecture about?} \textbf{Access control} answers one question:
``is \emph{this} user allowed to do \emph{this} action on \emph{this} file?''. In
\textbf{DAC}, the \textbf{owner} of a file decides who can use it (like on your laptop or
on Facebook). We see how to store these permissions (matrix, ACLs, capabilities, roles,
groups), a classic trap (the \emph{confused deputy}), and how Linux and Windows do it in
practice. \\
\textbf{Why is it useful?} Access control is the \textbf{first line of defense} and is
everywhere: apps, databases, operating systems, hardware. Understanding it lets you read
and set file permissions correctly (\texttt{ls -l}, \texttt{chmod}, \texttt{sudo}) and
spot programs that can be tricked into using their rights for an attacker.
\end{rembox}

\section{What is access control?}

\begin{rembox}[title=Access control in the physical world]
We already know it: keys for homes and cars, the CAMIPRO card for EPFL rooms (a physical door
with digital support). Computers need the same thing for files, sockets, database rows, memory.
\end{rembox}

\begin{defnbox}[title=Access control]
Security mechanism that ensures that ``all accesses and actions on objects by principals
are \textbf{within} the security policy''.
\end{defnbox}

\begin{exbox}[title=Questions an access control system answers]
Can Alice read file \texttt{/users/Bob/readme.txt}? Can Bob open a TCP socket to
\texttt{http://www.abc.com/}? Can Charlie write to row 15 of the table GRADES? \\
The answer is either \textbf{authorized} (``has permission'') or
\textbf{unauthorized} (``access denied''): only events within the security policy may
happen.
\end{exbox}

\begin{rembox}[title=Why is access control so important?]
Access control is the \textbf{first line of defense}, thus \textbf{it is used
everywhere}:
\begin{itemize}
  \item \textbf{Applications} --- online social networks, email servers, cloud storage;
  \item \textbf{Middleware} --- database management systems (DBMS);
  \item \textbf{Operating system} --- access to files, directories, ports, \ldots;
  \item \textbf{Hardware} --- memory, registers, privileges.
\end{itemize}
\end{rembox}

\begin{defnbox}[title=Authentication vs.\ authorization]
\begin{itemize}
  \item \textbf{Authentication} --- the system needs to \textbf{bind the actor to a
        principal} before authorization. A \textbf{principal} is an abstract entity
        that is authorized to act (users, connections, processes).
  \item \textbf{Authorization} --- the system needs to decide whether the
        \textbf{principal is authorized}: this is \textbf{access control}.
\end{itemize}
\end{defnbox}

\begin{warnbox}[title=Trusted Computing Base (TCB)]
It contains the mechanisms that do authentication and authorization: if either of them fails,
the security policy is breached, so both must be in the TCB (Lecture 2).
\end{warnbox}

\begin{rembox}[title={Origin of protection mechanisms (Lampson, 1971)}]
The original motivation was to keep \textbf{one user's malice or error from harming other
users}: by destroying/modifying another user's data, by reading/copying it without
permission, or by degrading the service another user gets (the extreme case being a
crash of the system). All these reasons are just as strong if ``user'' is replaced by
``\textbf{program}''.
\end{rembox}

\begin{goatbox}
Access control is the guard at the barn door. For every request ``this goat wants to do this
thing to this bucket'', the guard says yes or no, by looking at the rules. Before that, someone
must check that the goat really is that goat (authentication: ``show me your ear tag''); then the
guard decides (authorization: ``this goat may drink from this bucket'').
\end{goatbox}

\section{Implementing access control}

\begin{propbox}[title={``Checks soup'' vs.\ reference monitor}]
\begin{itemize}
  \item \textbf{What NOT to do: ``checks soup''} --- add checks all over the program,
        implementing the decision \textbf{in-line} based on the policy
        (\texttt{if user == Alice or user == Bob: open(\ldots)}). It is \textbf{error-prone}
        (a missed check? how to verify or debug?) and \textbf{hard to change} (a policy change
        means finding every check in thousands of lines).
  \item \textbf{What you SHOULD do: systematic calls to a ``reference monitor''} ---
        all over the program add checks that \textbf{call the monitor}: they state the
        authorization required and provide evidence about principals and objects; a
        \textbf{``central'' subsystem} establishes whether the checks pass or not
        (e.g.\ Apache Shiro: \texttt{subject.isPermitted("user:delete:jsmith")}).
\end{itemize}
\end{propbox}

\begin{rembox}[title=Least common mechanism? A tradeoff between principles]
The lecture's view: a central reference monitor \textbf{violates least common mechanism} (it is
shared by everybody; if it fails or is compromised, the policy can be breached), but it
\textbf{fulfils economy of mechanism} (access control in one place, easy to check) and
\textbf{complete mediation} (it can check every action). Not all principles can be followed at
once: here the one that is better for security wins. (One can also note that least common
mechanism is mostly about the isolation of \emph{data}, and a single well-tested monitor shares
code, not data.)
\end{rembox}

\begin{goatbox}
Do not give every farmhand a scribbled list of rules to check in their own way (``checks soup''):
someone will forget one. Put \emph{one} guard at the barn door who knows the rule book; everyone
asks the guard. If the rule changes, you tell one guard, not fifty farmhands.
\end{goatbox}

\section{Discretionary Access Control (DAC)}

\begin{defnbox}[title=Discretionary Access Control (DAC)]
\textbf{Object owners} decide how their resources are shared: the owner assigns the
permissions on what they own (e.g.\ Linux \texttt{chmod}, Windows, macOS, social
networks). DAC policies are often conceptualized as an \textbf{access control matrix}.
\end{defnbox}

\begin{rembox}[title=DAC in practice]
There is no central authority: objects are owned by different users. Since systems have huge
numbers of objects, owners usually choose among a \textbf{predefined set of options} rather than
fine-grained rules (``public / friends / only me''). Examples: operating systems, social networks
(Facebook, Instagram: the owner decides who sees a post), cloud file sharing (Google Drive,
Dropbox).
\end{rembox}

\subsection{The access control matrix}

\begin{defnbox}[title=Access control matrix]
An \textbf{abstract} representation of all \textbf{permitted} triplets
\textbf{(subject, object, access right)} within a system. DAC policies are often
conceptualized as an access control matrix (Lampson).
\begin{itemize}
  \item \textbf{Subjects (principals)} --- entities within an IT system (a user, a
        process, a service).
  \item \textbf{Objects (assets)} --- resources that (some) subject may access or use
        (a file, a folder, a row in a database, the system's memory, a machine in the
        network, a printer, a page in a website).
  \item \textbf{Operations} --- in abstract, subjects can \textbf{observe} and/or
        \textbf{alter} objects (read, write, append, execute).
\end{itemize}
\end{defnbox}

\begin{exbox}[title=Access control matrix]
$S = \{\text{Alice}, \text{Bob}\}$, $O = \{\text{file1}, \text{file2},
\text{file3}\}$, $A = \{\text{read}, \text{write}\}$.
\begin{center}
\begin{tabular}{l|c|c|c}
      & file1      & file2      & file3      \\ \hline
Alice & read, write &            & read       \\
Bob   &            & read, write & read, write
\end{tabular}
\end{center}
Can Alice read file1? \textbf{Yes.} Can Bob write file1? \textbf{No} (Bob has nothing on
file1). Can file3 be written by Alice? \textbf{No} (only read). \\
The matrix is more general than users and files (Lampson's original example was already richer):
an entry can be any operation (read, write, execute, replace, delete, \ldots).
\end{exbox}

\begin{warnbox}[title=The matrix is an abstract concept]
Not for direct implementation: size $O(f \cdot u)$ ($f$ files, $u$ users), which is huge
with thousands of files or hundreds of users. Moreover it is:
\begin{itemize}
  \item \textbf{usually very sparse} --- wasteful;
  \item \textbf{inefficient to access} --- need to keep the matrix in fast memory?
  \item \textbf{not extensible} --- add a file? add a user?
\end{itemize}
\end{warnbox}

\begin{goatbox}
Imagine a giant table: one row per goat, one column per bucket, and in each box what that goat
may do with that bucket. It answers every question, but with a million goats and a million
buckets it is a trillion boxes, almost all empty. So nobody stores the table: you either cut it
into columns (each bucket keeps its own list: ACL) or into rows (each goat keeps its own list:
capabilities).
\end{goatbox}

\subsection{Access Control Lists (ACLs)}

\begin{defnbox}[title=Access Control List (ACL)]
Associate permissions to \textbf{objects}: each object stores the \textbf{column} of the
matrix, i.e.\ the list of (subject, rights) pairs. \textbf{Blanks are not stored}.
\begin{center}
\texttt{file3: \{(Alice, read), (Bob, read/write)\}}
\end{center}
\end{defnbox}

\begin{propbox}[title=ACLs: pros and cons]
\textbf{Pros}: can be stored close to / with the resource; easy to determine who can
access a resource; easy to revoke rights by resource. \\
\textbf{Cons}: difficult to check at runtime (it is often unclear \emph{who} the subject
accessing is: see the confused deputy); difficult to audit all rights of a user (go file by file:
linear in the number of files); difficult to remove all permissions from a user (error-prone:
better remove their authentication credentials!); difficult delegation (hard to tell whether a
user may delegate, and to remove a temporary permission later).
\end{propbox}

\begin{goatbox}
ACL = a list stuck on each bucket: ``Alice may drink, Bob may drink and refill''. To know who
uses this bucket: read its list. To know everything Alice may do: walk to every bucket in the
farm.
\end{goatbox}

\subsection{Role Based Access Control (RBAC)}

\begin{defnbox}[title=Role Based Access Control (RBAC)]
Systems have too many subjects that come and go (large dynamic ACLs), but subjects are
similar to each other and need the same rights (a doctor has the same privileges as
another doctor).
\begin{enumerate}
  \item assign permissions to \textbf{roles};
  \item assign \textbf{roles} to subjects;
  \item subjects select an \textbf{active role} --- they have the permissions of the
        active role.
\end{enumerate}
Example: Alice the surgeon has both the Doctor and the Nurse roles. Acting as Nurse, she can read
a patient's treatment but not write it; switching to Doctor, she can do both. Acting in one role
at a time prevents misuse of privileges (\textbf{least privilege}).
\end{defnbox}

\begin{warnbox}[title=RBAC problems]
\begin{enumerate}
  \item \textbf{Role explosion} --- temptation to create fine-grained roles, denying the
        benefits of RBAC.
  \item \textbf{Limited expressiveness} of simple RBAC --- problems implementing least
        privilege; some roles are relative (``Alice's doctor'' vs.\ ``any doctor'').
  \item \textbf{Separation of duty is difficult} --- ``two doctors are needed to
        authorize a procedure'': the mechanism needs to ensure they are distinct!
\end{enumerate}
With too many roles (almost one per user) the system is as complex as an ACL; with few roles it
cannot express nuances (a doctor should only see \emph{their own} patients); and since all
subjects of a role look the same to the system, it cannot tell one doctor from two.
\end{warnbox}

\begin{goatbox}
Instead of listing goats one by one, give hats: ``milker hat'', ``shepherd hat''. The rules say
what each hat may do, and each goat gets some hats. A goat wears one hat at a time and has only
that hat's rights. Trouble starts when you invent a new hat for every goat (no gain), or when a
rule needs ``two \emph{different} shepherds'' (two shepherd hats look the same).
\end{goatbox}

\subsection{Group Based Access Control}

\begin{defnbox}[title=Group Based Access Control]
Observation: some permissions are always needed together (e.g.\ access to sockets and
network interface).
\begin{enumerate}
  \item assign permissions to access objects to \textbf{groups};
  \item assign subjects to \textbf{groups};
  \item subjects have the permissions of \textbf{all} their groups.
\end{enumerate}
\textbf{Negative permissions} (explicit denials for a subject) can be added to implement
fine-grained policies. Example: group1 gives access to file1, file2, file3; Alice needs file2 and
file3, so group1 makes sense for her, but the policy forbids her file1. Instead of a new group
(or breaching the policy), add a negative permission ``Alice: no read/write on file1''.
\end{defnbox}

\begin{propbox}[title=Negative permissions are checked first]
If group permissions were checked first, a group would grant the access and the
negative permission would never be applied. Checking \textbf{negative permissions first}
is the fail-safe choice (deny overrides allow).
\end{propbox}

\begin{rembox}[title=RBAC vs.\ groups]
With \textbf{roles}, a subject has the permissions of the \emph{one active} role; with
\textbf{groups}, a subject has the union of the permissions of \emph{all} its groups.
\end{rembox}

\begin{goatbox}
Groups are clubs: the ``water club'' always gets the bucket \emph{and} the tap together. A goat in
three clubs has everything of the three clubs at once. If one club member must not touch one
bucket, write a red ``NO'' card for her, and the guard reads red cards first.
\end{goatbox}

\subsection{Capabilities}

\begin{defnbox}[title=Capabilities]
Associate permissions to \textbf{subjects}: each subject stores the \textbf{row} of the
matrix, i.e.\ the list of (object, rights) pairs. \textbf{Blanks are not stored}.
\begin{center}
\texttt{Alice: \{(file1, read/write), (file3, read)\}}
\end{center}
\end{defnbox}

\begin{propbox}[title=Capabilities: pros and cons]
\textbf{Pros}: can be stored with the subject (portable!); easy to audit all subject
permissions; delegating is ``simple''. \\
\textbf{Cons}: revoking just one permission is hard (if the objects are unstructured, finding the permission
is linear in the number of objects of the user); transferability --- once the capability is given, how can we
prevent sharing? (if Alice gives Bob her right to read file1, which is not in the matrix, the
policy is broken); authenticity --- how to check it? (like checking a concert ticket is not
fake). \\
Portability is a real advantage in \textbf{distributed systems}: the subject shows its credential
to any node, wherever the object is. The paper ``Capability myths demolished'' argues that the
cons are addressable.
\end{propbox}

\begin{keybox}[title=ACL vs.\ capabilities]
\begin{center}
\begin{tabular}{l|l|l}
                       & \textbf{ACL}           & \textbf{Capabilities} \\ \hline
Stored with            & the object (column)    & the subject (row)     \\
Who can access $o$?    & easy                   & hard                  \\
What can $s$ access?   & hard                   & easy                  \\
Revocation             & easy per object        & hard per object       \\
Delegation             & difficult              & ``simple''
\end{tabular}
\end{center}
\end{keybox}

\begin{rembox}[title=History of capability-based security]
Introduced by Dennis \& Van Horn (1966); used in the Plessey System 250 (1972), CAP
computer (1976), Intel iAPX 432 (1981). The industry then adopted UNIX and its ACLs
(1982, Linux 1991). Revival: seL4 formally verified microkernel (2009), CHERI (2015),
Google Fuchsia (2016), WASI (2019), ARM Morello (2022), Scala capabilities at EPFL.
\end{rembox}

\begin{goatbox}
Capabilities = each goat carries its own keyring: ``key to bucket 1 (drink, refill), key to
bucket 3 (drink)''. Everything Alice may do is on her keyring; she can lend a key easily. But to
take back every key of bucket 3, you must find every goat; and a goat can copy or fake a key unless
keys are hard to forge.
\end{goatbox}

\subsection{Ambient authority and the confused deputy}

\begin{defnbox}[title=Ambient authority]
\textbf{Ambient authority} is used by a subject if, for an action to succeed, it
\textbf{only} needs to specify the \textbf{operation} and the \textbf{names} of the
involved object(s). The \textbf{subject} (with authority) is \textbf{implicit}:
in \texttt{open("file1", "rw")} the subject is missing, it is understood to be the
process owner.
\end{defnbox}

\begin{propbox}[title=Ambient authority: pros and cons]
\textbf{Pro}: no need to repeat the subject all the time (usability). \\
\textbf{Cons}: the program cannot check permissions for whom it acts; least privilege
becomes harder to enforce; \textbf{confused deputy problem}!
\end{propbox}

\begin{rembox}[title={ACLs, capabilities and ambient authority}]
ACL systems generally use ambient authority: permissions are checked for the user \emph{running}
the program. With capabilities, the capability itself carries the right (and who holds it):
there is \textbf{no ambient authority}.
\end{rembox}

\begin{defnbox}[title=Confused deputy problem]
Problem with ambient authority: a \textbf{privileged program} can be \textbf{tricked to
misuse its authority}. A confused deputy is a privileged principal used by a less privileged
principal to increase her access to the system.
\end{defnbox}

\begin{exbox}[title=The CAMIPRO story]
Students want to enter a classroom their CAMIPRO cards do not open. Breaking the lock or EPFL's
central services is hard. Easier: ask a new professor, who does not know well who may enter
where, to open the door with his very privileged badge. The professor is the confused deputy.
\end{exbox}

\begin{exbox}[title=Pay-per-use compiler]
The compiler receives \texttt{(input, output)}, compiles \texttt{input}, writes usage
into \texttt{bill} and errors into \texttt{output}.
\begin{center}
\begin{tabular}{l|c|c|c}
         & input & output      & bill        \\ \hline
Alice    & write & read        & read        \\
Compiler & read  & read, write & read, write
\end{tabular}
\end{center}
Alice cannot write \texttt{bill}, but she calls the compiler with
\texttt{(input, bill)}: the compiler writes the errors into \texttt{output = bill}
\textbf{with its own authority} $\Rightarrow$ \texttt{bill} is corrupted. The compiler
is the \textbf{confused deputy}. (Same idea: a student with a card asks someone
authorized to badge them into a room.)
\end{exbox}

\begin{algbox}[title=How to avoid confused deputies]
Ambient authority is used for convenience in many real systems (OS services, web
servers, \ldots). Solutions:
\begin{enumerate}
  \item re-implement access control in the privileged process (removes the ambient authority,
        but expensive and cumbersome, thus error-prone);
  \item let the privileged process check authorization for Alice (the ambient authority stays,
        but it checks that the user invoking it has the same permissions);
  \item \textbf{capabilities can help}! The compiler holds a capability for
        \texttt{bill}; to compile, Alice must give a capability for the
        \texttt{output} file --- she cannot give a capability for writing
        \texttt{bill}, so she cannot confuse anyone.
\end{enumerate}
\end{algbox}

\begin{goatbox}
The sheepdog may enter every pen. A sly goat says ``dog, please push the gate of \emph{that}
pen'', and the dog, who only asks ``may \emph{I} push it?'', obeys: the goat used the dog's
rights, not her own. That is the confused deputy. The cure: the dog asks ``may \emph{you}?'', or
the goat must hand the dog a key to the pen --- a key she does not have.
\end{goatbox}

\section{DAC in practice: UNIX}

\subsection{Principals and architecture}

\begin{defnbox}[title=UNIX principals and groups]
\begin{itemize}
  \item \textbf{User Identities (UIDs)} and \textbf{Group Identities (GIDs)}:
        originally 16-bit (now 32-bit) numbers. Special UIDs: $-2$, $0$, $1$, \ldots
  \item Each user has its own directory \texttt{/home/username}; user accounts are in
        \texttt{/etc/passwd} (\texttt{username:password:UID:GID:info:home:shell}: \texttt{info}
        is a comment, \texttt{home} the directory where the user lands at login, \texttt{shell}
        the default shell).
  \item Users belong to one or more groups: primary group (\texttt{/etc/passwd}),
        other groups (\texttt{/etc/group}).
\end{itemize}
\end{defnbox}

\begin{propbox}[title=UNIX security architecture]
\begin{itemize}
  \item \textbf{Everything is a file}; each user ``owns'' a set of files.
  \item Each file has a simple \textbf{ACL} expressing its access control policy. The
        owner can \textbf{modify the permissions of the file} (but cannot give away
        ownership). System files are owned by special users that can make system
        operations.
  \item All processes run by a user run \textbf{with that user's privileges}
        $\Rightarrow$ \textbf{ambient authority}!
\end{itemize}
\end{propbox}

\subsection{File permissions}

\begin{defnbox}[title=UNIX file ACLs]
Each file is assigned an \textbf{owner UID} and a \textbf{GID}, and has \textbf{9
permission bits}: 3 actions (\textbf{r}ead, \textbf{w}rite, e\textbf{x}ecute) $\times$ 3
subjects (\textbf{owner}, \textbf{group}, \textbf{other}). Plus 3 attributes:
\texttt{suid}, \texttt{sgid} and \texttt{sticky}. \\
Different semantics for \textbf{directories}: Read $\to$ list files, Write $\to$ add
files (create, move or rename a file in it), Exec $\to$ traverse the directory (\texttt{cd}).
In practice one cannot read or write the files of a directory without \texttt{x} on it. \\
\textbf{other} = anyone who is neither the owner nor in the file's group.
\end{defnbox}

\begin{exbox}[title=Reading \texttt{ls -l}]
\texttt{drwxrwxr-x 1 catronco catronco 4096 Sep 16 14:23 exampledir} \\
First char: \texttt{d} = directory, \texttt{-} = file; then \texttt{rwx} for owner,
\texttt{rwx} for group, \texttt{r-x} for others; then number of links, owner, group,
size, last modified, filename. \\
Other lines of the slide: \texttt{-rwxrwxrw- \ldots hello} (a file, owner and group may run it,
others read and write), \texttt{-rw-rw--w- \ldots test1.txt} (others may write but not
read). \\
The owner changes permissions with \texttt{chmod}: \texttt{+r}, \texttt{-w},
octal \texttt{666}, \texttt{662}, sticky \texttt{+t} or \texttt{1666}, suid
\texttt{+s} or \texttt{4666}.
\end{exbox}

\begin{algbox}[title=UNIX access control in action]
Compare the \texttt{UID}/\texttt{GID} of the \textbf{process} trying to perform the
action with the state of the file (owner, group, mode bits). \textbf{Order matters}:
\begin{enumerate}
  \item if the \texttt{UID} says you are the owner: check the bits for \textbf{owner};
  \item if not owner, but your group is the file's group: check the bits for
        \textbf{group};
  \item otherwise check the bits for \textbf{other}.
\end{enumerate}
The \texttt{root} user is \textbf{never denied access}.
\end{algbox}

\begin{warnbox}[title=Order matters]
Only \emph{one} class of bits is checked: an owner with \texttt{---} for owner but
\texttt{rwx} for group is denied, even if they are in the group.
\end{warnbox}

\begin{goatbox}
Every UNIX file has a tiny label with three boxes: what its owner may do, what its group may do,
what everybody else may do (each box: read, write, run). The guard looks at the goat: owner? use
the first box. Not owner but in the group? the second. Otherwise the third. Only one box is ever
read. And root, the farmer, walks through every door.
\end{goatbox}

\subsection{Special users and bits}

\begin{defnbox}[title=Super user (\texttt{root})]
Special user account with \textbf{User ID 0}: accesses system files and special
operations, can access anything ((almost) all security checks disabled).
\textbf{\texttt{root} is in the TCB!} \\
\textbf{Never login as root} (some distributions give root no password): use
\texttt{sudo} (run \emph{one} command as root) or \texttt{su} (switch to another user, root by
default --- dangerous, since nothing you do afterwards is checked). \texttt{sudo su catronco}:
use your sudo rights to become user \texttt{catronco} without knowing their password. Normal users
still need system services that run with system privileges: hence \texttt{suid}.
\end{defnbox}

\begin{defnbox}[title=\texttt{suid} / \texttt{sgid}]
Setuid and setgid bits indicate that a file is \textbf{not run with the privileges of
the launcher}, but \textbf{with the privileges of the owner} user/group. Useful for
normal users to access system services that need system privileges, respecting least
privilege, e.g.\ changing a password:
\texttt{-rwsr-xr-x 1 root root \ldots /bin/passwd}. When the program ends, the privileges go back
to normal.
\end{defnbox}

\begin{warnbox}[title=Setuid root programs are dangerous]
They are in the TCB: how do you know that a \texttt{suid} program does what it is meant
to do, and \emph{only} what it is meant to do? (e.g.\ a setuid-root \texttt{python}
gives root to anyone.)
\end{warnbox}

\begin{exbox}[title=Why a setuid-root program is dangerous]
\begin{itemize}
  \item It runs as \textbf{root whoever launches it}, and the launcher controls its
        inputs (arguments, environment, file names). A single bug (buffer overflow, path
        trick, command injection) gives the attacker \textbf{root}: a
        \textbf{confused deputy}.
  \item A \textbf{general-purpose} program must never be setuid root: with
        \texttt{-rwsr-xr-x root python}, anyone can run
        \texttt{python -c 'import os; os.setuid(0); os.system("/bin/sh")'} and get a
        root shell. (The \texttt{python2.7} of the slide is \texttt{-rwxr-xr-x}, no
        \texttt{s}: it is \emph{not} setuid, it shows what to avoid.)
  \item It breaks \textbf{least privilege}: full root power for one small action.
\end{itemize}
\textbf{Mitigation}: few setuid-root programs, small and audited (economy of mechanism,
like \texttt{/bin/passwd}); drop privileges once the privileged step is done; prefer
finer-grained rights (Linux capabilities, \texttt{sudo} limited to one command); never
make an interpreter or an editor setuid.
\end{exbox}

\begin{defnbox}[title=Sticky bit]
``Restricted deletion bit'' (\texttt{chmod +t}).
\begin{itemize}
  \item \textbf{Directories}: prevents unprivileged users from removing or renaming a
        file in the directory \textbf{unless they own the file} (e.g.\ \texttt{/tmp}:
        users can only edit their own files).
  \item \textbf{Files}: historically prevented a program from being moved from swap
        for fast load; Linux currently ignores the bit.
\end{itemize}
On \texttt{/tmp} (shared by all users, e.g.\ for downloads) the sticky bit ensures that only a
file's owner (or root) can delete or rename it.
\end{defnbox}

\begin{defnbox}[title=Nobody]
Special user (\textbf{User ID $-2$}) that owns no files and belongs to no user. Safer
user to execute code you do not know (particularly obfuscated code): limits damages if it
misbehaves / gets compromised. It is a way to \textbf{sandbox} programs: even if the program is
malicious and takes control of its user, it gains no privileges. (Further reading: ``Setuid
demystified'', 2002; ``Jails: confining the omnipotent root'', 2000.)
\end{defnbox}

\begin{goatbox}
\texttt{suid}: a program that wears its owner's hat while it runs, whoever starts it. The
password-change program wears root's hat so that any goat can change \emph{its own} password in
the farmer's book. Dangerous if the program can be tricked into writing anything else in the
book. Sticky bit: in the shared shed (\texttt{/tmp}), you may only throw away your own things.
Nobody: a goat with no name and no belongings, to run strangers' programs safely.
\end{goatbox}

\subsection{Exercise}

\begin{exbox}[title={Exercise --- ACL in Linux}]
\texttt{drwxr-x--- 2 manager\_A core\_devs 4096 Sep 24 11:35 /project\_beta} \\
\textbf{manager\_A} is the directory owner, a member of \texttt{core\_devs}, and is
sudo; \textbf{developer\_B} is a member of \texttt{core\_devs}; \textbf{intern\_C} is
\textbf{not} a member of \texttt{core\_devs}. \\
For each user, determine which of \texttt{cd /project\_beta}, \texttt{ls /project\_beta}
and \texttt{touch /project\_beta/<file>} succeed or fail, and explain why based on the DAC
rules. \\[4pt]
Same state, with a new subdirectory and a file inside it: \\
\texttt{drwxr-xr-x 2 developer\_B core\_devs 4096 Sep 24 12:00 /project\_beta/Foo} \\
\texttt{-rw-r--r-- 1 developer\_B core\_devs 512 Sep 24 12:05 \ldots/Foo/intern\_project.txt}
\begin{enumerate}[label=Q\arabic*:]
  \item \texttt{/project\_beta} is owned by manager\_A, but \texttt{Foo} is owned by
        developer\_B. Find a scenario (sequence of commands) that would have led to this
        situation.
  \item intern\_C runs \texttt{cd /project\_beta/Foo} then
        \texttt{cat intern\_project.txt}. Do they work?
  \item intern\_C runs \texttt{cat /project\_beta/Foo/intern\_project.txt} from their
        home directory. Does it work?
  \item Who do they need to contact to fix the problem?
\end{enumerate}
\end{exbox}

\begin{algbox}[title={Solution --- ACL in Linux}]
\texttt{/project\_beta} is \texttt{drwxr-x---}, i.e.\ owner \texttt{rwx}, group
\texttt{r-x}, other \texttt{---}. Only \emph{one} class is
checked (owner, else group, else other). On a directory: \texttt{r} = list (\texttt{ls}),
\texttt{w} = create/delete files (\texttt{touch} of a new file), \texttt{x} = traverse
(\texttt{cd}).
\begin{center}
\begin{tabular}{l|l|c|c|c}
 & class used & \texttt{cd} (\texttt{x}) & \texttt{ls} (\texttt{r}) & \texttt{touch} (\texttt{w}) \\ \hline
manager\_A   & owner: \texttt{rwx}   & OK     & OK     & OK     \\
developer\_B & group: \texttt{r-x}   & OK     & OK     & denied \\
intern\_C    & other: \texttt{---}   & denied & denied & denied
\end{tabular}
\end{center}
\begin{itemize}
  \item \textbf{manager\_A} is the owner: the owner bits apply (being sudo is not
        needed).
  \item \textbf{developer\_B} is not the owner but is in \texttt{core\_devs}: the group
        bits apply; no \texttt{w} $\Rightarrow$ cannot create \texttt{notes.txt}.
  \item \textbf{intern\_C} falls in ``other'' (\texttt{---}): everything is denied.
\end{itemize}
\textbf{Q1:} developer\_B cannot create files in \texttt{/project\_beta} (group has no
\texttt{w}). The lecture's answer: a normal, deliberate situation created by a superuser (root,
or manager\_A with sudo), who created the directory and transferred its ownership:
\texttt{sudo mkdir /project\_beta/Foo; sudo chown developer\_B /project\_beta/Foo} --- delegating
part of the project without giving away the parent directory. (Another scenario: manager\_A runs
\texttt{chmod g+w /project\_beta}, developer\_B runs \texttt{mkdir Foo} and creates the file, then
manager\_A runs \texttt{chmod g-w /project\_beta}.) \\
\textbf{Q2:} no. \texttt{cd /project\_beta/Foo} fails: intern\_C has no \texttt{x} on
\texttt{/project\_beta} (``other'' is \texttt{---}). \texttt{Foo} being \texttt{r-x} for
others does not help: \textbf{every directory on the path} must be traversable. The
\texttt{cat} then runs from the unchanged working directory and does not find the file. \\
\textbf{Q3:} no, same reason: \texttt{Permission denied} on \texttt{/project\_beta}, even
though the file itself is world-readable (\texttt{r--}). \\
\textbf{Q4:} \textbf{manager\_A}, the owner of \texttt{/project\_beta} (and sudo):
developer\_B owns \texttt{Foo} but not the parent directory. manager\_A can add intern\_C
to \texttt{core\_devs} (least-privilege-wise better than \texttt{chmod o+x}, which would
open the directory to everybody).
\end{algbox}

\begin{rembox}[title={Open question of the slides: what about links? (answer not in the slides)}]
A \textbf{symbolic link} is only a path: following it, the system checks the permissions of the
target \emph{and} \texttt{x} on every directory of the target's path, so a symlink placed in
intern\_C's home does not help. A \textbf{hard link} is another name for the same file (same
inode, same permission bits): if one existed in a directory intern\_C can traverse, intern\_C
could read the world-readable file without crossing \texttt{/project\_beta}. That is why modern
Linux only lets you create a hard link to a file you own or can read and write
(\texttt{protected\_hardlinks}).
\end{rembox}

\section{DAC in practice: Windows}

\begin{defnbox}[title=Windows access control]
Principals = users, machines, groups, \ldots; objects = files, registry keys, printers,
\ldots
\begin{itemize}
  \item Each object has a \textbf{Discretionary Access Control List (DACL)}.
  \item Each process (or thread) has an \textbf{access token} with: the login user
        account (the process ``runs as'' this user), all groups of which the user is a
        member (recursively!), all privileges assigned to these groups.
\end{itemize}
The DACL is compared with the process' access token \textbf{when creating a handle to the
object}: if there is no permission, no handle is given.
\end{defnbox}

\begin{defnbox}[title=DACL]
A list of \textbf{Access Control Entries (ACEs)}, each with a \textbf{type}
(negative/positive, i.e.\ access denied/allowed), a \textbf{principal}, and
\textbf{permissions} (more fine-grained than UNIX), plus flags and others. Negative
entries come \textbf{first} so that a denial cannot be bypassed by an earlier allow.
\end{defnbox}

\begin{exbox}[title=DACL evaluation]
ACE1: \textit{deny} Andrew read/write/execute; ACE2: \textit{allow} Group A write; ACE3:
\textit{allow} Everyone read/execute. Thread A (Andrew, groups A, B, C) is denied by ACE1.
Thread B (Jane, group A) is granted write (ACE2) and read/execute (ACE3).
\end{exbox}

\begin{rembox}[title=Least privilege by default]
Windows runs with least privilege by default; elevated rights require
``Run as administrator''.
\end{rembox}

\begin{goatbox}
Windows does the same as UNIX with longer labels: each object has a list of lines ``allow / deny,
who, what'' (deny lines first, so a ``no'' is never skipped), and each running program carries a
badge with its user, all its groups and their privileges. The guard compares badge and list once,
when the program asks to open the object.
\end{goatbox}

\section{Mandatory Access Control (MAC)}

\begin{defnbox}[title=Mandatory Access Control (MAC)]
A \textbf{central security policy} assigns the permissions, and \textbf{users cannot
override it}, even for what they created. Used by organizations with a need for central
controls.
\end{defnbox}

\begin{exbox}[title=Where is MAC used?]
\begin{itemize}
  \item \textbf{Military} --- focus on \textbf{confidentiality} (classification levels:
        a document is readable only by users with a high enough clearance, whatever its
        owner wants).
  \item \textbf{Hospital} --- focus on \textbf{confidentiality and integrity} (records
        unchanged, and seen only by those who need them to treat the patient). The
        hospital owns the patient records and limits their sharing; regulatory
        requirements (e.g.\ \textbf{HIPAA} for health information) may limit it further.
  \item \textbf{Banking} --- focus on \textbf{integrity} (e.g.\ whoever enters a transfer
        cannot approve it).
\end{itemize}
In the military, integrity matters too but less: better that battle plans be corrupted or
deleted than leaked to the enemy. In a bank, leaking someone's earnings is bad press, but
modifying the records to move money is critical. \textbf{Modern access control systems are a
hybrid of DAC and MAC.}
\end{exbox}

\section{DAC vs.\ MAC}

\begin{center}
\begin{tabularx}{\linewidth}{@{}l X X@{}}
\toprule
 & \textbf{DAC} & \textbf{MAC} \\
\midrule
Who sets the permissions & The \textbf{owner} of the object & A \textbf{central authority} (the policy) \\
Can a user change them?  & Yes, for what they own         & No, even for what they created \\
Typical use & Windows, Linux, macOS, social networks & Military, hospitals, banking \\
\bottomrule
\end{tabularx}
\end{center}

\begin{rembox}[title=It is not a question of scale]
What separates them is \textbf{who has the authority to decide}, not the size of the
system: a huge system can be DAC (a social network), a small one MAC. Test: if Alice
creates a file, can she decide by herself who may read it? Yes $\Rightarrow$ DAC; no,
the policy decides $\Rightarrow$ MAC.
\end{rembox}

\begin{goatbox}
DAC: each goat decides who may use \emph{its} things. MAC: the farm has one rule book, and no goat
can change it, not even for things it made itself.
\end{goatbox}

\section{Summary}

\begin{keybox}[title=Summary]
\begin{itemize}
  \item \textbf{Discretionary Access Control}: owners establish permissions (in
        \textbf{MAC}, a central policy does).
  \item Conceptualized as an \textbf{access control matrix}.
  \item Implemented in two flavors: \textbf{access control lists} (permissions
        associated to objects) and \textbf{capabilities} (permissions associated to
        subjects).
  \item When relying on access control, always think about \textbf{confused deputies}.
\end{itemize}
\end{keybox}

\begin{recapbox}[title={Exam-ready recap --- Lecture 3: Access control basics and DAC}]
\textbf{Access control} = the mechanism ensuring every (principal, object, action) is within the
policy; answer: authorized / denied. First line of defence, everywhere (apps, DBMS, OS,
hardware). \textbf{Authentication} binds the actor to a principal, \textbf{authorization} (access
control) decides; both in the TCB.
\begin{itemize}
  \item \textbf{Implementation}: no ``checks soup''; calls to a central \textbf{reference monitor}
        (economy of mechanism + complete mediation, at the cost of least common mechanism).
  \item \textbf{DAC}: the \textbf{owner} decides. \textbf{MAC}: a \textbf{central policy}
        decides, users cannot override (military: C; hospital: C+I; bank: I). Test: can the
        creator decide alone who reads? Not a question of scale. Real systems: hybrids.
  \item \textbf{Access control matrix} (Lampson): subjects $\times$ objects $\to$ rights.
        Abstract only: $O(f \cdot u)$, sparse, slow, not extensible.
  \item \textbf{ACL} = column, stored with the object: easy ``who can access $o$'', easy revoke
        per object; hard ``what can $s$ do'', hard to remove a user, hard delegation.
        \textbf{Capabilities} = row, stored with the subject: portable, easy audit and
        delegation; hard revocation, transferability, authenticity.
  \item \textbf{Groups}: subject gets the \emph{union} of all its groups;
        \textbf{negative permissions checked first} (fail-safe). \textbf{RBAC}: permissions
        $\to$ roles $\to$ subjects, \emph{one active role} (least privilege); problems: role
        explosion, limited expressiveness, separation of duty.
  \item \textbf{Ambient authority}: the subject is implicit (\texttt{open("f","rw")} runs as the
        process owner) $\Rightarrow$ \textbf{confused deputy}: a privileged program tricked into
        using its rights for someone else (pay-per-use compiler writing errors into
        \texttt{bill}; professor's badge). Fixes: AC inside the deputy, deputy checks the
        caller, \textbf{capabilities} (no ambient authority).
  \item \textbf{UNIX}: UIDs/GIDs, \texttt{/etc/passwd}, everything is a file, owner sets
        permissions (\texttt{chmod}) but cannot give ownership; 9 bits rwx $\times$
        owner/group/other; on directories r = \texttt{ls}, w = create, x = \texttt{cd} (needed
        on \emph{every} directory of the path). \textbf{Only one class checked}: owner, else
        group, else other; root (UID 0, in the TCB) never denied. \texttt{suid/sgid}: run with
        the \emph{owner's} privileges (\texttt{passwd}; dangerous, in the TCB); sticky bit: only
        the owner deletes/renames in the directory (\texttt{/tmp}); nobody (UID $-2$): sandbox.
  \item \textbf{Windows}: DACL of ACEs (deny first) per object, access token per process (user,
        groups recursively, privileges), checked when creating a handle; least privilege by
        default.
\end{itemize}
\textbf{How it builds on what came before}: access control is the first \emph{mechanism} (Lecture
1) that enforces a \emph{policy} on principals and assets, and it is the principles of Lecture 2
in action: the reference monitor is \textbf{complete mediation} (with a small TCB:
\textbf{economy of mechanism}), negative permissions first and whitelists are \textbf{fail-safe
defaults}, active roles and need-to-know are \textbf{least privilege}, two distinct doctors is
\textbf{separation of privilege}, setuid programs and the confused deputy show what happens when
a privileged component is in the TCB. Next lecture: MAC, where the policy itself is designed with
models.
\end{recapbox}
